\documentclass[10pt,a4paper]{article}
\usepackage{amsmath,amssymb,amsthm,graphicx}
% Minimal retained preamble for lane-p2 cards taken from
% sources/ahl-16-fixed-points-2019/editorial-preamble.tex
% (doi:10.5802/ahl.16, Annales Henri Lebesgue 2 (2019) 149-185, CC BY 4.0).
%
% Copied verbatim from the editorial preamble: the package list, the
% theorem declarations under the shared counter and the macro block, so that
% every macro and environment of the retained spans is available.
%
% Omitted and why: the class line (the assembler supplies the class); the
% a4paper geometry line and the fontspec block with its explicit lmroman10
% files (encodings and fonts fight the forced xelatex engine, and the
% assembler supplies the page setup).
% Added for this card only: the two enumerated items of the retained statement are quoted by
% the proof through the paper numbers 3.4.1 and 3.4.2, so the label of the top level of that
% enumeration is set globally to those numbers (enumerate redefines its label locally).
\makeatletter
\gdef\theenumi{3.4.\arabic{enumi}}
\makeatother
\usepackage{amsmath,amssymb,amsthm,mathtools,graphicx,xcolor,hyperref,xurl,xspace,ifthen}
\newtheorem{theo}{Theorem}[section]
\newtheorem{lemm}[theo]{Lemma}
\newtheorem{theorem}[theo]{Theorem}
\newtheorem*{remas*}{Remarks}
\providecommand{\og}{«\nobreakspace}
\providecommand{\fg}{\nobreakspace»}
\newtheorem{lem}[theo]{Lemma}
\newtheorem{lemma}[theo]{Lemma}
\newtheorem{prop}[theo]{Proposition}
\newtheorem{coro}[theo]{Corollary}
\newtheorem{corol}[theo]{Corollary}
\newtheorem{conj}[theo]{Conjecture}
\newtheorem{Proposition}[theo]{Proposition}
\newtheorem{theoreme}[theo]{Theorem}
\theoremstyle{definition}
\newtheorem{defi}[theo]{Definition}
\newtheorem{exam}[theo]{Example}
\newtheorem{exem}[theo]{Example}
\newtheorem{rema}[theo]{Remark}
\newtheorem{nota}[theo]{Notation}
\newtheorem*{rema*}{Remark}
\newtheorem*{defi*}{Definition}
\newtheorem*{theo*}{Theorem}
\newtheorem*{lemm*}{Lemma}
\newtheorem*{prop*}{Proposition}
\newcommand{\enoncename}{}
\newtheorem{corpusenonce}[theo]{\enoncename}
\newenvironment{enonce}[2][]{\renewcommand{\enoncename}{#2}\begin{corpusenonce}}{\end{corpusenonce}}
\newenvironment{enonce*}[2][]{\par\medskip\noindent\textbf{#2.}\itshape}{\par\medskip}
\newenvironment{proofc}[1][\proofname]{\begin{proof}[#1]}{\end{proof}}
\newcommand{\Mc}{\mathcal{M}}
\newcommand{\pointir}{.\ ---\ }
\providecommand{\mathof}[1]{\mathrm{#1}}

\numberwithin{equation}{section}
\numberwithin{figure}{section}
\numberwithin{table}{section}
\def\endappendix{}
%~Mouliné par MaN_auto v.0.21.4 2019-04-30 09:02:09


\usepackage[all]{xy}

\hyphenation{auto-morphism
every
Roelcke
peri-od-ic
peri-od-icy
rep-re-sen-ta-tion
}








\newcommand{\A}{\mathcal A}
\newcommand{\N}{\mathbb{N}}
\newcommand{\Z}{\mathbb{Z}}
\newcommand{\Q}{\mathbb{Q}}
\newcommand{\K}{\mathcal{K}}
\newcommand{\CC}{\mathbb{C}}

\newcommand*{\mA}{\mathbf{A}}
\newcommand*{\mB}{\mathbf{B}}
\newcommand*{\mC}{\mathbf{C}}
\newcommand*{\mD}{\mathbf{D}}
\newcommand*{\mE}{\mathbf{E}}
\newcommand*{\mF}{\mathbf{F}}
\newcommand*{\mH}{\mathbf{H}}
\newcommand*{\mZ}{\mathbf{Z}}
%\newcommand*{\mA}{\mathbf{A}}
%\newcommand*{\mA}{\mathbf{A}}

\newcommand{\restrict}[2]{#1 \upharpoonright #2}

\newcommand{\e}{\varepsilon}

\DeclareMathOperator{\Aut}{Aut}
\DeclareMathOperator{\Stab}{Stab}
\DeclareMathOperator{\Age}{Age}

\newcommand{\actson}{\curvearrowright}

\newcommand{\RUCB}{\mathrm{RUC}_{b}}
\newcommand{\ROB}{\mathrm{Ro}_{b}}
\newcommand{\WAP}{\mathrm{WAP}}
\newcommand{\AP}{\mathrm{AP}}

\graphicspath{{./figures/}}

\newcommand*{\mk}{\mkern -1mu}
\newcommand*{\Mk}{\mkern -2mu}
\newcommand*{\mK}{\mkern 1mu}
\newcommand*{\MK}{\mkern 2mu}

\hypersetup{urlcolor=purple, linkcolor=blue, citecolor=red}

\newcommand*{\romanenumi}{\renewcommand*{\theenumi}{\roman{enumi}}}
\newcommand*{\Romanenumi}{\renewcommand*{\theenumi}{\Roman{enumi}}}
\newcommand*{\alphenumi}{\renewcommand*{\theenumi}{\alph{enumi}}}
\newcommand*{\Alphenumi}{\renewcommand*{\theenumi}{\Alph{enumi}}}
\let\oldexists\exists
\renewcommand*{\exists}{\mathrel{\oldexists}}
\let\oldforall\forall
\renewcommand*{\forall}{\mathrel{\oldforall}}


\usepackage{cleveref}
\crefname{originalproposition}{original proposition}{original propositions}
\Crefname{originalproposition}{Original proposition}{Original propositions}
\makeatletter
\newlabel{prop:aOSRamsey}{{3.3}{source}{}{sourcepointer-1}{}}
\newlabel{prop:aOSRamsey@cref}{{[originalproposition][0][]3.3}{[1][0][]source}{}{}{}}
\makeatother
\title{Finite oscillation stability of an algebra is equivalent to the Ramsey property up to twice the precision}
\author{Lionel Nguyen Van Thé}
\date{doi:10.5802/ahl.16; CC BY 4.0}
\begin{document}
\maketitle
\setcounter{section}{3}
\setcounter{enumi}{3}
\setcounter{equation}{1}
\noindent\textit{Editorial note.} The counted unit is the article's own Proposition prop:OSRP at source lines 340-346 together with its proof at source lines 348-352, a straight adjacent pairing in the source file. The source is the publisher version of the article in Annales Henri Lebesgue 2 (2019), pages 149-185, DOI 10.5802/ahl.16, released under CC BY 4.0.\par\medskip
\noindent\textit{Editorial note.} The retained proof invokes the article's oscillation stability proposition in its Ramsey form, which is a statement this lane has already counted on another card of the same source. It is declared here as a navigation pointer carrying only its printed number, 3.3, and none of its text.\par\medskip
\noindent\textit{Editorial note.} The card sets the section counter to three and the label of the top level of the retained enumeration to 3.4, so that the two enumerated items display as 3.4.1 and 3.4.2, which are the numbers the article prints for them and the numbers the retained proof refers to by equation reference; the equation counter is set so the displayed formula of the proof keeps the number the article gives it. The definition of the item label lives in this card's own copy of the preamble, because the enumerate environment redefines its label locally.\par\medskip
\noindent\textit{Editorial note.} The licence evidence used here is the observation file kept with the source, license-observation.json, which records the licence as CC BY 4.0, points at the inherited publisher-version evidence and states that the landing page displays no licence line. No acceptance record is written anywhere.\par\medskip
\noindent\textit{Editorial note.} The retained preamble is this source's editorial preamble with three omissions only: the class line, which the assembler supplies; the page geometry; and the fontspec block with its explicit font files, which would fight the forced xelatex engine. Everything else is retained verbatim; the card copies the shared slice kept at focused-seeds/logic/P2-ahl16-retained-preamble.tex and adds the item label definition named above.\par\medskip
\noindent\textit{Editorial note.} Neither the retained statement nor the retained proof contains a citation, so this card needs no bibliography entry. The gate probe on the idle interval reported the interval free and returned RESULT=OK.\par\medskip
\noindent\textit{Allowed standard prerequisites (editorial).}\begin{itemize}
\item Finitely oscillation stable families and the Ramsey property up to a precision. A finite family of functions is finitely oscillation stable when for every finite set of group elements and every precision some translate makes all the functions constant up to that precision; the article also says that a structure has the Ramsey property up to a precision for colorings in a family when for every finite substructure there is a copy of it on which every coloring of the family is constant up to that precision. The author's own definitions are the ones used by the counted statement and proof.. The counted proposition is the quantitative bridge between the analytic condition on families of functions and the combinatorial condition on colorings.
\item The oscillation stability proposition in its Ramsey form. The article's earlier proposition states that for a Fraisse structure and a unital left-invariant closed star-subalgebra, a finite set of finite colorings is finitely oscillation stable exactly when for every finite substructure there is a copy on which every coloring of the set is constant. The author's own statement is the one used by the counted proof.. The forward direction of the counted proof invokes that proposition to turn the constancy up to a precision obtained from the algebra into the Ramsey-type conclusion.
\end{itemize}\medskip
\ifdefined\hypertarget\hypertarget{sourcepointer-1}{}\fi
\noindent\textit{Source reference.} Original proposition 3.3: sources/ahl-16-fixed-points-2019/body-01.tex, source lines 318--325. Navigation pointer invoked by the retained proof at the step that turns the constancy of a finite set of colorings up to a precision into the existence of a copy of any finite substructure. That proposition is an earlier result of the article and is not retained in this card, so the reference is declared as a pointer carrying the number the article itself prints for that statement, that is the section together with the place of the statement inside it, read off an enumeration of every numbered theorem-like environment of the source in order: it is 3.3; the cited statement is an allowed input and none of it is reproduced inside the span.\par\medskip
\begin{prop}\label{prop:OSRP}
Let $\mF$ be a Fra\"{\i}ss\'e structure, and let $\A$ be a unital, left-invariant, closed C$^{*}$-subalgebra of $\RUCB(\Aut(\mF))$. TFAE:
\begin{enumerate}
\item \label{3.4.1} Every finite $\mathcal F\subset \A$ is finitely oscillation stable.
\item \label{3.4.2} For every $\e>0$, $\mF$ has the Ramsey property up to $2\e$ for colorings in $(\A)_{\e}$.
\end{enumerate}
\end{prop}

\begin{proof}
Assume that every finite $\mathcal F\subset \A$ is finitely oscillation stable. Fix $\mA$ in $\Age(\mF)$, $\mathcal C \subset (\A)_{\e}$ a finite set of finite colorings of $\binom{\mF}{\mA}$, $H\subset \Aut(\mF)$ finite. Fix $\{ f_{\chi}: \chi\in \mathcal C\}\subset \A$ and $\eta >0$ so that $\| \chi - f_{\chi}\|_{\infty}+\eta<\e$ for every $\chi \in \mathcal C$. By finite oscillation stability of $\{ f_{\chi}: \chi\in \mathcal C\}$, find $g\in \Aut(\mF)$ so that every $f_{\chi}$ is constant up to $\eta$ on $Hg$. Then every $\chi \in \mathcal C$ is constant up to $2\e$ on $Hg$. Thanks to Proposition~\ref{prop:aOSRamsey}, we deduce that for every $\mB \in \Age(\mF)$, there exists $b\in \binom{\mF}{\mB}$ such that every $\chi\in \mathcal C$ is constant up to $2\e$ on $\binom{b(\mB)}{\mA}$. This is exactly what we needed to prove.

Conversely, assume that~\eqref{3.4.2} holds, and fix $\mathcal F\subset \mathcal A$ finite, $\e>0$, $H\subset \Aut(\mF)$ finite. Let $\{ \chi_{f} : f\in \mathcal F\}$ be a finite family of finite colorings in $(\A)_{\e/4}$ so that $\| f-\chi_{f}\|_{\infty}<\e/4$ for every $f\in \mathcal F$. Thanks to Proposition~\ref{prop:aOSRamsey}, \eqref{3.4.2}~implies that there is $g\in \Aut(\mF)$ so that every $\chi_{f}$ is constant up to $\e/2$ on $Hg$. Then, every $f\in \mathcal F$ is constant up to $\e$ on $Hg$ and $\mathcal F$ is finitely oscillation stable.
\end{proof}

\end{document}
